\begin{textsample}{POS Dim 3 – vem\_ed\_11 – Score 4.00 – vem\_ed\_11/t049.txt}  \label{ex:f3_pos_139}
BOAS PRÁTICAS Pedro Sérgio Rosa é doutor em Educação para a Ciência (Unesp), mestre em Física (Unicamp), Bacharel (UEL) e Licenciado (Unesp) em Física.
Professor do curso de Produção Fonográfica na Fatec Tatuí, relata como superar a barreira do idioma e desenvolver um Projeto Colaborativo Internacional em uma área aparentemente árida, a física das ondas.
Os desafios são vários, mas, o principal é, em \textbf{tese}, a barreira do idioma.
Por que em \textbf{tese}?
Porque são três aspectos: leitura, escrita e comunicação \textbf{oral}.
A leitura para mim era normal, pois, desde o Bacharelado em Física estava acostumado com os textos científicos em inglês.
Por outro lado, a escrita em outro idioma não é nada trivial, uma vez que é necessário estar continuamente estudando.
A parte mais complexa é a comunicação \textbf{oral}: compreender o que um nativo da língua fala e responder adequadamente.
Classifico por nível de dificuldade, a partir do mais fácil: leitura, escrita e comunicação \textbf{oral} em língua inglesa.
Para um nativo, a comunicação entre pares é sempre muito mais rápida.
Em relação ao ensino de Acústica Aplicada à Produção Fonográfica, disciplina que ministro na Fatec, observo que, para planejar um Projeto Colaborativo Internacional (PCI), são necessárias muitas reuniões, diálogos e correções de rota.
O que me ajudou enormemente foi ter realizado, com a professora de inglês Dulce Helena Soares Villa Nova e com o estudante Flavio Lee, da monitoria em línguas da Fatec Tatuí, aulas e correções na pronúncia, bem como treinos que tornaram minha audição mais aguda.
Isso permitiu, com o tempo, melhorar minha compreensão da fala de Machele Kindle e James Goodwin, professores colaboradores no PCI desenvolvido entre Fatec Tatuí e BridgeValley Community College (EUA).
Porém, como todo processo de aprimoramento, para se ter proficiência em outra língua, a quantidade de horas de prática é decisiva.
O contato com o nativo cria uma memória de audibilidade que permite assistir a um programa, filme, documentário na língua inglesa sem a necessidade da leitura total da tradução pela legenda.
Não há que se ter receio de aprender algo novo ou de adquirir novas habilidades.
A dica é não ter medo de errar, mas, se ocorrer o erro, pedir que o professor nativo ou o professor de línguas que participa do PCI ajude.
Durante o projeto, passei a assistir a filmes com áudio e legenda em inglês, para treinar a audição: creio que isso ajudou também.
Finalizando, diria que os Projetos Colaborativos Internacionais da Cesu / Centro Paula Souza precisam ser explorados pela comunidade de alunos e professores que consideram a \textbf{cooperação} como um ganho fundamental para seu aprimoramento, tanto pessoal quanto profissional.

% matched lemmas: cooperação, oral, tese
\end{textsample}
